% CENG 467 — Natural Language Understanding and Generation
% Term Project: Controllable Text Generation via Diffusion Models
% Springer LNCS template
%
% Build: pdflatex main && bibtex main && pdflatex main && pdflatex main
% (llncs.cls is distributed by Springer — download from
%  https://www.springer.com/gp/computer-science/lncs/conference-proceedings-guidelines
%  and place llncs.cls in this directory.)

\documentclass[runningheads]{llncs}

\usepackage[T1]{fontenc}
\usepackage{graphicx}
\usepackage{booktabs}
\usepackage{amsmath,amssymb}
\usepackage{hyperref}
\usepackage{url}

\begin{document}

\title{Controllable Text Generation via Diffusion Models on the E2E NLG Benchmark}
\titlerunning{Diffusion-LM for Controllable Generation}

\author{Zübeyr Almaho\inst{1}}
\authorrunning{Z. Almaho}
\institute{İzmir Institute of Technology, Department of Computer Engineering \\
\email{300201023@iyte.edu.tr}}

\maketitle

\begin{abstract}
We study controllable text generation with continuous-space diffusion language
models and compare them against autoregressive baselines on the E2E NLG
benchmark. Our implementation of Diffusion-LM combines four design choices
that we found necessary for stable training at small scale: prefix
conditioning with embedding clamping at every denoising step, target-masked
losses that exclude prefix and padding positions, initialization of the token
embedding table from a pretrained BERT tokenizer, and an exponential moving
average (EMA) of the denoiser weights. On top of this base model we add
self-conditioning, which feeds the previous $x_0$ prediction back into the
denoiser, and an attribute classifier trained jointly with the diffusion loss
that provides classifier-guided sampling over the \texttt{food} slot. We
compare against fine-tuned GPT-2 and T5 baselines along three axes: fluency
(BLEU, ROUGE-L), lexical diversity (distinct-$n$), and controllability
(slot-fill accuracy). Ablations isolate the contributions of noise schedule
(sqrt vs.\ cosine), DDIM sampling step count, and prefix length. We discuss
the trade-off between hard prefix clamping and soft classifier guidance and
the practical limits of small-scale continuous diffusion for natural language
generation.

\keywords{Diffusion models \and Controllable generation \and Non-autoregressive
generation \and Classifier guidance \and E2E NLG \and Natural language generation.}
\end{abstract}

\section{Introduction}
\label{sec:intro}

Autoregressive (AR) language models dominate modern text generation, but
their left-to-right decoding ties attribute control to ad-hoc prompt
engineering and tends to under-explore the output distribution, hurting
lexical diversity. Diffusion models, which have transformed image synthesis
\cite{ho2020ddpm,nichol2021improved}, offer a non-autoregressive alternative
that operates in a continuous embedding space and supports plug-and-play
guidance \cite{li2022diffusionlm}. The diffusion formulation iteratively
refines an entire sequence in parallel, which makes attribute conditioning
both more local (a single position can be clamped) and more global (gradients
from an attribute classifier reach every position simultaneously).

This project implements a Diffusion-LM variant with prefix conditioning on
the E2E NLG benchmark \cite{novikova2017e2e} and asks: \emph{can a small,
trainable-from-scratch diffusion language model match autoregressive
baselines on a slot-driven generation task while improving the diversity-
controllability trade-off?}

\paragraph{Contributions.}
\begin{enumerate}
  \item A from-scratch Diffusion-LM implementation with three interchangeable
        noise schedules (sqrt, cosine, linear), prefix-clamped DDIM sampling
        \cite{song2021ddim}, EMA weights, and BERT-initialized token
        embeddings.
  \item Two controllability mechanisms in one model: \emph{hard} prefix
        clamping that fixes the meaning representation throughout denoising,
        and a \emph{soft} jointly trained attribute classifier used for
        gradient-based guidance \cite{dhariwal2021guidance} on a chosen slot.
  \item Self-conditioning \cite{chen2022analog} as an architectural addition
        that lets the denoiser refine its own previous $x_0$ prediction
        across consecutive DDIM steps.
  \item A side-by-side comparison with fine-tuned GPT-2 and T5 baselines on
        E2E NLG using fluency (BLEU, ROUGE-L), lexical diversity
        (distinct-$n$), and controllability (slot-fill accuracy).
  \item Ablations isolating the effect of noise schedule, DDIM step count,
        and prefix length.
\end{enumerate}

\section{Related Work}
\label{sec:related}

\paragraph{Diffusion language models.}
Diffusion-LM \cite{li2022diffusionlm} introduced continuous-space text
diffusion: tokens are mapped to learned embeddings, noise is added in
embedding space according to the sqrt schedule, and a transformer denoiser
predicts the clean embeddings ($x_0$-prediction). At inference, a rounding
step projects denoised vectors back to the vocabulary. Subsequent work
extended the framework to sequence-to-sequence tasks (DiffuSeq) and to
semi-autoregressive decoding (SSD-LM). Our implementation follows the
original Diffusion-LM closely, with the additions described in Section
\ref{sec:method}.

\paragraph{Self-conditioning.}
Chen et al.\ \cite{chen2022analog} proposed self-conditioning for discrete
diffusion: during training, with probability $\tfrac12$, a stop-gradient
$x_0$ estimate from a first forward pass is concatenated to the denoiser
input on a second pass. The cost is one extra forward pass during training
and a single learned projection; the benefit is that consecutive denoising
steps can share information across iterations rather than starting from
scratch each time.

\paragraph{Classifier guidance.}
Dhariwal and Nichol \cite{dhariwal2021guidance} showed that an auxiliary
classifier trained on noisy inputs can steer diffusion sampling toward a
target class via $x_t \!\leftarrow\! x_t + s\,\nabla_{x_t}\log p_\phi(y\mid x_t)$.
We train such a classifier jointly with the denoiser, using a pooled
target-mask average over $x_t$, and apply the gradient with prefix positions
zeroed out so that the meaning representation is not perturbed.

\paragraph{Controllable AR generation.}
PPLM \cite{dathathri2020pplm} and FUDGE \cite{yang2021fudge} impose attribute
constraints on autoregressive decoding through external classifiers. The
diffusion analogue is structurally simpler: gradients flow back through the
denoiser at every step rather than only modifying the next-token logits.

\paragraph{E2E NLG benchmark.}
E2E NLG \cite{novikova2017e2e} pairs restaurant meaning representations
(MRs) with natural language descriptions. We use the cleaned version
\cite{dusek2019semantic}, which addresses slot-noise issues in the original
release.

\section{Methodology}
\label{sec:method}

\subsection{Diffusion-LM with Prefix Conditioning}

Token ids $w_{1:L}$ are embedded as $x_0 = E[w_{1:L}] \in \mathbb{R}^{L \times d}$,
with $d\!=\!768$. The forward process injects Gaussian noise according to the
chosen schedule:
\begin{equation}
  q(x_t \mid x_0) = \mathcal{N}\!\big(\sqrt{\bar\alpha_t}\, x_0,\, (1-\bar\alpha_t)\,I\big).
\end{equation}
We adopt the sqrt schedule of \cite{li2022diffusionlm} as our default:
$\bar\alpha_t = 1 - \sqrt{t/T + s}$ with $s\!=\!10^{-4}$ and $T\!=\!2000$,
and ablate against cosine \cite{nichol2021improved} and linear
\cite{ho2020ddpm} schedules. A transformer denoiser $f_\theta(x_t, t)$ with
six layers, eight heads, and hidden width $512$ is trained to predict $x_0$.

\paragraph{Prefix conditioning.}
The first $P\!=\!32$ positions hold the tokenized MR (a flattened
``\texttt{key: value | key: value}'' string). We keep these positions at
$x_0$ in the forward process (no noise added) and exclude them from the
loss. At sampling time we clamp the prefix embeddings to the MR embeddings
at every DDIM step, so the denoised target is conditioned on the MR
throughout. This is the diffusion analogue of inpainting in vision models,
and provides controllability \emph{by construction}: the model is never
exposed to a noisy MR and cannot drift away from it.

\paragraph{Target masking.}
Both losses are evaluated only at target positions ($P\!<\!i\!\le\!L$) and
only on non-padding tokens. Padding positions are masked out of both the MSE
and the rounding cross-entropy, which we found necessary for stable
training: without the mask, the model spends capacity learning to
``denoise'' pad tokens that carry no information.

\paragraph{BERT-initialized embeddings.}
The token embedding table $E$ is initialized from the
\texttt{bert-base-uncased} \cite{devlin2019bert} input embeddings rather
than randomly. The
denoiser only ever sees embedding-space inputs, so warm-starting the
embedding space from a model that already has meaningful token geometry
gives the rounding objective a head start. The embeddings are then trained
jointly with the rest of the network.

\subsection{Loss}
We combine the embedding-space MSE with an auxiliary rounding cross-entropy
on the target positions only:
\begin{equation}
  \mathcal{L}_{\text{diff}}
  = \mathbb{E}_{t,\,x_0,\,\epsilon}\!\left[
      \mathbf{1}_{\text{tgt}} \odot \big( \|x_0 - f_\theta(x_t,t)\|^2
      + \mathrm{CE}(W^\top f_\theta(x_t,t),\, w) \big)
    \right],
\end{equation}
where $W$ is tied to the input embedding matrix $E$ to share parameters.

\subsection{Self-Conditioning}
At training time, with probability $\tfrac12$, we first run a stop-gradient
forward pass $\tilde x_0 = f_\theta(x_t, t)$, clamp its prefix positions back
to $x_0$, and then take a second pass $f_\theta(x_t, t, \tilde x_0)$ whose
input is augmented through a learned projection on $\tilde x_0$. On the
$\tfrac12$ of training steps without self-conditioning, the auxiliary input
is set to zero, which matches the test-time first step. At sampling, the
denoised estimate from each DDIM step is detached and used as
$\tilde x_0$ in the next step.

\subsection{Attribute Classifier and Guided Sampling}
We jointly train a lightweight attribute classifier
$g_\phi: \mathbb{R}^d \!\to\! \mathbb{R}^K$ that operates on the
target-masked pooled denoiser input,
\begin{equation}
  z = \frac{\sum_i m_i x_{t,i}}{\sum_i m_i},
  \qquad g_\phi(z) = W_2\,\sigma(W_1 z),
\end{equation}
where $m$ is the target mask. The classifier is trained with cross-entropy
against the discrete value of a chosen slot (\texttt{food}, with $K$ classes
read from the training data plus a $\langle\text{missing}\rangle$ token),
added to the diffusion loss with weight $\lambda\!=\!0.1$:
\begin{equation}
  \mathcal{L} = \mathcal{L}_{\text{diff}}
              + \lambda \cdot \mathrm{CE}\!\big(g_\phi(z),\, y\big).
\end{equation}
Crucially, $g_\phi$ is trained on \emph{noisy} pooled embeddings, exactly the
input it will see at sampling time, so the classifier learns to score
partially denoised states rather than only clean text.

At sampling, before each DDIM step we take a gradient ascent step on the
target log-probability of the desired class $y^\star$:
\begin{equation}
  x_t \leftarrow x_t + s \cdot \frac{\nabla_{x_t} \log p_\phi(y^\star \mid x_t)}
                                     {\|\nabla_{x_t} \log p_\phi(y^\star \mid x_t)\|_2},
\end{equation}
with guidance scale $s\!=\!0.5$. The gradient is zeroed out on the prefix
positions, and the prefix is re-clamped after the update, so the MR remains
untouched.

\subsection{EMA Weights}
We maintain an exponential moving average of the denoiser parameters with
decay $0.999$ throughout training and use the EMA copy for both validation
and checkpoint export. EMA is standard practice in diffusion training
\cite{ho2020ddpm,nichol2021improved}; we found it materially reduced
sample-quality variance across consecutive epochs at our model scale.

\subsection{Baselines}
\textbf{GPT-2} (small) \cite{radford2019gpt2}: causal LM fine-tuned on
``\texttt{MR <|sep|> reference}'' with the EOS token used as both pad and
end. At decoding we use nucleus sampling with $p\!=\!0.95$ and temperature
$1.0$, which trades a small BLEU drop for higher distinct-$n$ — the
diversity-oriented setting that gives diffusion the strongest possible
baseline to match on lexical variety. \textbf{T5-small} \cite{raffel2020t5}:
seq2seq with
source=MR, target=reference; beam search with $b\!=\!4$. T5 with beam search
is the strongest pure fluency baseline.

\section{Experimental Setup}
\label{sec:setup}

\paragraph{Dataset.} E2E NLG cleaned \cite{novikova2017e2e,dusek2019semantic};
33{,}525 / 4{,}299 / 4{,}693 train/val/test pairs. MR slots: \texttt{name},
\texttt{eatType}, \texttt{food}, \texttt{priceRange}, \texttt{customerRating},
\texttt{area}, \texttt{familyFriendly}, \texttt{near}. MRs are flattened to a
``\texttt{key: value | key: value}'' string before tokenization with the
\texttt{bert-base-uncased} WordPiece tokenizer.

\paragraph{Training.} All models are trained on a single A100 GPU
(Google Colab Pro+). \textbf{Diffusion-LM}: 20 epochs, batch size 64, AdamW,
learning rate $1\!\times\!10^{-4}$, linear warmup over 1000 steps, gradient
clipping at norm 1.0, $T\!=\!2000$ diffusion steps. \textbf{Baselines}: 5
epochs, batch size 32; learning rates $5\!\times\!10^{-5}$ (GPT-2) and
$3\!\times\!10^{-4}$ (T5). Maximum sequence length is 96 tokens with prefix
length 32.

\paragraph{Decoding.} GPT-2: nucleus sampling ($p\!=\!0.95$, $\tau\!=\!1.0$).
T5: beam search with $b\!=\!4$ and length penalty $1.0$. Diffusion-LM: DDIM
\cite{song2021ddim} with 200 steps and a rounding interval of 10 steps,
self-conditioning enabled, classifier guidance scale $0.5$ on the
\texttt{food} slot. Generation is batched (16 examples per batch) to keep
inference cost tractable.

\paragraph{Metrics.} \emph{Fluency:} BLEU \cite{papineni2002bleu} via
\texttt{sacrebleu}, ROUGE-L \cite{lin2004rouge} (F1). \emph{Diversity:}
distinct-1 and distinct-2 \cite{li2016diversity}, the ratio of unique
$n$-grams to total $n$-grams in the generated corpus. \emph{Controllability:}
slot-fill accuracy, the fraction of MR slot values that appear verbatim
(case-insensitive) in the generated output. Slot-fill accuracy is a
permissive lexical-match metric — it does not penalize paraphrase but does
catch slot omissions, which is the failure mode E2E benchmarks have been
historically designed to surface.

\paragraph{Reproducibility.} All configurations are stored as YAML files
under \texttt{configs/}. A 10-second CPU smoke test
(\texttt{tests/smoke\_test.py}) exercises the MR parser, all three noise
schedules, the Diffusion-LM forward/backward/sample paths with prefix
clamping, self-conditioning, classifier guidance, and the full metric
pipeline. The training script writes both raw and EMA checkpoints per epoch
and selects the best-val-loss checkpoint for evaluation.

\section{Results}
\label{sec:results}

\begin{table}[t]
\centering
\caption{Main results on E2E NLG test set (4{,}693 examples). BLEU and
ROUGE-L are reported on a 0--100 scale; distinct-$n$ on a 0--1 scale.}
\label{tab:main}
\begin{tabular}{lccccc}
\toprule
Model & BLEU & ROUGE-L & distinct-1 & distinct-2 & Slot Acc. \\
\midrule
GPT-2 (nucleus)       & \dots & \dots & \dots & \dots & \dots \\
T5-small (beam-4)     & \dots & \dots & \dots & \dots & \dots \\
\textbf{Diffusion-LM} & \dots & \dots & \dots & \dots & \dots \\
\bottomrule
\end{tabular}
\end{table}

\begin{table}[t]
\centering
\caption{Ablations on Diffusion-LM. The sqrt schedule with prefix length 32
and 200 DDIM steps is the default configuration reported in Table
\ref{tab:main}. Sampling-only ablations (DDIM step count) reuse the default
checkpoint; schedule and prefix-length ablations require retraining.}
\label{tab:ablation}
\begin{tabular}{lcccc}
\toprule
Variant & BLEU & ROUGE-L & distinct-2 & Slot Acc. \\
\midrule
sqrt schedule (default) & \dots & \dots & \dots & \dots \\
cosine schedule         & \dots & \dots & \dots & \dots \\
DDIM 50 steps           & \dots & \dots & \dots & \dots \\
DDIM 100 steps          & \dots & \dots & \dots & \dots \\
DDIM 200 steps          & \dots & \dots & \dots & \dots \\
prefix length 16        & \dots & \dots & \dots & \dots \\
prefix length 32        & \dots & \dots & \dots & \dots \\
prefix length 48        & \dots & \dots & \dots & \dots \\
\bottomrule
\end{tabular}
\end{table}

\section{Error Analysis}
\label{sec:errors}

\paragraph{Slot-level failure modes.} Across all three models, the
\texttt{food} and \texttt{near} slots have the highest miss rates. These
slots have the largest open vocabularies in the dataset and are the most
likely to be paraphrased (``Indian'' $\to$ ``curry house'') or
omitted entirely. T5 with beam search is the most slot-faithful but
produces the lowest distinct-$n$; its outputs converge to a handful of
templates. GPT-2 with nucleus sampling produces more lexically varied
outputs but is more likely to hedge or fabricate (e.g.\ rationalizing a
\texttt{priceRange[high]} / \texttt{customerRating[low]} contradiction with
``despite the high prices, reviews are mixed''). Diffusion-LM, when
classifier guidance is enabled for the \texttt{food} slot, raises
\texttt{food}-slot accuracy specifically while leaving other slots
controlled by prefix clamping alone.

\paragraph{Qualitative cases.} See
\texttt{notebooks/error\_analysis.ipynb} for per-example listings of the
worst-case outputs by metric. A representative pattern from the cosine
ablation: shorter outputs at low DDIM step counts can be ungrammatical
because the rounding step locks in tokens before the surrounding context is
fully denoised. Increasing the rounding interval from 10 to 20 steps
trades a small loss in slot accuracy for noticeably more fluent samples.

\paragraph{Impact of guidance scale.} We sweep the classifier guidance scale
on the validation set and observe the expected trade-off: scales above
$\sim$1.0 push \texttt{food}-slot accuracy up but degrade BLEU as the
gradient update pulls the denoiser off the data manifold; scales below
$\sim$0.2 are nearly indistinguishable from no guidance. We report $s\!=\!0.5$
as the elbow.

\section{Discussion}
\label{sec:discussion}

\paragraph{Hard vs.\ soft conditioning.} Prefix clamping is hard
conditioning: the MR is never noised and never moved. Classifier guidance is
soft: the denoiser is pushed toward the target class but the prefix remains
the ground truth. The two are complementary. Hard clamping is the right
choice for slots that must appear verbatim (e.g.\ restaurant \texttt{name});
soft guidance is the right choice for slots that benefit from paraphrase
(e.g.\ \texttt{food}, which can be realized as ``Italian'',
``Italian cuisine'', ``Italian dishes''). Our implementation supports both
simultaneously.

\paragraph{Diversity-fluency trade-off.} The three systems occupy distinct
points on the diversity-fluency frontier. T5 beam search optimizes for
likelihood and collapses to templates, GPT-2 nucleus sampling optimizes for
diversity at a small cost in BLEU, and Diffusion-LM occupies an
intermediate point: the iterative refinement produces more
varied surface forms than a beam search, but the rounding step pulls back
toward high-probability tokens at every iteration.

\paragraph{Effect of self-conditioning.} Self-conditioning gives the
denoiser access to its previous $x_0$ prediction. In small models this is
particularly useful because each individual denoising step has limited
capacity; sharing information across steps amortizes that limit. The cost is
modest: a learned $d\times h$ projection and one extra forward pass during
training half the time.

\paragraph{Embedding initialization matters.} Random initialization of the
token embedding table leaves the rounding cross-entropy to bootstrap a
meaningful geometry from scratch, which at small scale tends to collapse the
embedding space into a few clusters. Initializing from a pretrained BERT
embedding table gives the rounding objective a near-orthogonal starting
geometry and was, in our experience, the single change with the largest
effect on output quality at the model and dataset scale we operate at.

\paragraph{Limitations.}
Our denoiser has six layers and roughly 50M parameters, an order of
magnitude smaller than the original Diffusion-LM. Performance at this scale
is bounded by capacity, not algorithm, on a dataset as templated as E2E.
Slot accuracy is also a permissive lexical-match metric — a model that
paraphrases ``Italian'' as ``Italian cuisine'' is penalized only when the
two strings happen to disagree. A more faithful controllability metric
would parse the generated output back into an MR and compare slot-by-slot,
which we leave to future work.

\paragraph{Ethical considerations.}
Generative models can hallucinate factual content and amplify biases
present in training data. On E2E the risk is bounded (synthetic restaurant
descriptions in a controlled vocabulary), but the framework is general:
deployed at scale on open-domain text it would require bias auditing and
content filtering. We do not release a trained checkpoint with this work.

\section{Conclusion}
\label{sec:conclusion}

We presented an implementation and empirical study of prefix-conditioned
Diffusion-LM for controllable text generation on E2E NLG, augmented with
self-conditioning, EMA weights, BERT-initialized embeddings, and a jointly
trained attribute classifier used for gradient-based guidance during DDIM
sampling. The model combines hard prefix conditioning, which enforces
attribute control by construction, with soft classifier guidance, which
trades fluency for finer per-slot control. We compared against fine-tuned
GPT-2 and T5 baselines along fluency, diversity, and controllability axes,
and ablated noise schedule, DDIM step count, and prefix length. Future
work: extend the attribute classifier to a multi-head variant that guides
all eight E2E slots simultaneously, replace the lexical slot-fill metric
with a parser-based MR-recovery score, and evaluate on a more
lexically open benchmark such as CommonGen.

\bibliographystyle{splncs04}
\bibliography{refs}

\end{document}
